% Source-derived chapter 05: Spectral realization of Weyl algebras
% Original source: tates-thesis-de-rham-setting
\chapter{Spectral realization of Weyl algebras}
\label{tates-thesis-de-rham-setting-chapter-05}
\source{tates-thesis-de-rham-setting}

\begin{remark}[Source section]
\label{tates-thesis-de-rham-setting-chapter-05-source}
\source{tates-thesis-de-rham-setting}
\source{tates-thesis-de-rham-setting}
This chapter reproduces the corresponding section of the retrieved
LaTeX source, including its definitions, statements, examples, and
proofs. It has not yet been translated into Lean.
\notready
\end{remark}

% The source section title was: Spectral realization of Weyl algebras
\label{s:weyl}
\source{tates-thesis-de-rham-setting}

\subsection{Overview}

\subsubsection{}

Let $i_n:\sZ^{\leq n} \into \sY_{\log}^{\leq n}$ denote the 
canonical embedding, as in \S \ref{ss:y-not}.

We have $i_{n,*}^{\IndCoh}(\sO_{\sZ^{\leq n}}) \in \Coh(\sY_{\log}^{\leq n})$,
which lies in $\IndCoh^*(\sY_{\log}^{\leq n})^{\heart}$ by
Proposition \ref{p:zn-cl}.

Let $\sF_n \coloneqq 
\Av_!^{\Gr_{\bG_m}^{\leq n},w}
(i_{n,*}(\sO_{\sZ^{\leq n}}))$, which is a compact object in $\IndCoh^*(\sY^{\leq n})$.

In this section, we construct an action of 
a Weyl algebra in $2n$ generators on $\sF_n$ (considered as an object of 
$\IndCoh^*(\sY^{\leq n})$).

\subsubsection{}

To formulate our construction more canonically, let
$W_n$ denote the algebra of global differential operators
on the scheme $\fL_n^+ \bA^1 = \Spec\Sym(t^{-n}k[[t]]dt/k[[t]]dt)$. 

Let $W_n^{op}$ denote $W_n$ with the reversed multiplication.
We will construct a canonical homomorphism:
%
\begin{equation}\label{eq:action}
\source{tates-thesis-de-rham-setting}
W_n^{op} \to \ul{\End}_{\IndCoh^*(\sY^{\leq n})}(\sF_n) \in \Alg = \Alg(\Vect).
\end{equation}

\begin{rem}
\source{tates-thesis-de-rham-setting}
\notready

In Proposition \ref{p:ff}, we will show that this map is actually an isomorphism. 

\end{rem}

\subsection{Reduction to generators and relations}\label{ss:coconn-reduction}
\source{tates-thesis-de-rham-setting}

Note that $\sF_n$ lies in the heart of the $t$-structure constructed
in Proposition \ref{p:y-t-str} (or rather, its truncated counterpart, as in 
\S \ref{ss:y-trun-indcoh}).

Therefore, the right hand side of \eqref{eq:action} 
lies in cohomological degrees $\geq 0$.

As the left hand side of \eqref{eq:action} is certainly in cohomological 
degree $0$, it suffices to construct a homomorphism:
%
\[
W_n^{op} \to \tau^{\leq 0} \ul{\End}_{\IndCoh^*(\sY^{\leq n})}(\sF_n) = 
H^0\ul{\End}_{\IndCoh^*(\sY^{\leq n})}(\sF_n) \in \Alg(\Vect^{\heart}).
\]

As both terms are now in degree $0$ and the left hand side
has a standard algebra presentation, such a construction 
may be given by constructing generators and checking relations.

We need to construct maps:
%
\begin{equation}\label{eq:fns}
\source{tates-thesis-de-rham-setting}
\begin{gathered}
t^{-n}k[[t]]dt/k[[t]]dt = (k[[t]]/t^n)^{\vee} \to 
\ul{\End}_{\IndCoh^*(\sY^{\leq n})}(\sF_n) \\
(\omega \in t^{-n}k[[t]]dt/k[[t]]dt) \mapsto \vph_{\omega}
\end{gathered}
\end{equation}

\noindent and:
%
\begin{equation}\label{eq:vfs}
\source{tates-thesis-de-rham-setting}
\begin{gathered}
k[[t]]/t^n \to 
\ul{\End}_{\IndCoh^*(\sY^{\leq n})}(\sF_n) \\
(f \in k[[t]]/t^n) \mapsto \xi_{f}.
\end{gathered}
\end{equation}

\noindent We then need to check that the $\vph_{\omega}$ operators
mutually commute, that the $\xi_f$ operators mutually commute, and
the identity:\footnote{The sign on the right hand side reflects working
with $W_n^{op}$, not $W_n$.}

\begin{equation}\label{eq:uncertainty}
\source{tates-thesis-de-rham-setting}
[\xi_f,\vph_{\omega}] = -\Res(f\omega) \cdot \id \in  
H^0\ul{\End}_{\IndCoh^*(\sY^{\leq n})}(\sF_n).
\end{equation}

We provide these constructions and check these identities in what follows.

\subsection{Action of functions}\label{ss:fn-intro}
\source{tates-thesis-de-rham-setting}

First, we construct the map $\omega \mapsto \vph_{\omega}$ 
from \eqref{eq:fns}.

We assume the reader is familiar with the notation
from \S \ref{s:aj} below. 

\subsubsection{}\label{ss:fn-concl}
\source{tates-thesis-de-rham-setting}

We have a commutative diagram:
%
\[
\xymatrix{
\Gr_{\bG_m}^{pos,\leq n} \times \sZ^{\leq n} \ar[d]^{\alpha}
 \ar[dr]^{p_1} 
& \\
\Gr_{\bG_m}^{pos,\leq n} \times \sZ^{\leq n} \ar[r]^{p_1}
& \Gr_{\bG_m}^{pos,\leq n}
}
\]

\noindent where $\alpha = (p_1,\act)$ for
$\act:\Gr_{\bG_m}^{pos,\leq n} \times \sZ^{\leq n} \to \sZ^{\leq n}$
the action morphism from above.

We have an evident isomorphism:\footnote{One can work with
$\IndCoh^*$ or $\QCoh$ for our purposes here. 
We chose
$\IndCoh^*$ for the sake of definiteness. 

We also have
written $(-)^*$ in place of $(-)^{*,\IndCoh}$ to keep
the notation simpler. To be clear, this notation refers to
the left adjoint to the $\IndCoh^*$ pushforward functor.}
%
\[
\omega_{\Gr_{\bG_m}^{pos,\leq n}} \boxtimes 
\sO_{\sZ^{\leq n}} = 
p_1^*(\omega_{\Gr_{\bG_m}^{pos,\leq n}}) \simeq
\alpha^*p_1^*(\omega_{\Gr_{\bG_m}^{pos,\leq n}}) \in 
\IndCoh^*(\Gr_{\bG_m}^{pos,\leq n} \times \sZ^{\leq n})
\]

\noindent giving a morphism:
%
\[
\omega_{\Gr_{\bG_m}^{pos,\leq n}} \boxtimes 
\sO_{\sZ^{\leq n}} \to 
\alpha_*^{\IndCoh}(\omega_{\Gr_{\bG_m}^{pos,\leq n}} \boxtimes 
\sO_{\sZ^{\leq n}}).
\]

\noindent Applying $p_{2,*}^{\IndCoh}$, we obtain a canonical
morphism:
%
\[
\Gamma^{\IndCoh}(\omega_{\Gr_{\bG_m}^{pos,\leq n}}) \otimes  
\sO_{\sZ^{\leq n}} \to 
p_{2,*}^{\IndCoh}
\alpha_*^{\IndCoh}(\omega_{\Gr_{\bG_m}^{pos,\leq n}} \boxtimes 
\sO_{\sZ^{\leq n}}) = 
\act_*^{\IndCoh}(\omega_{\Gr_{\bG_m}^{pos,\leq n}} \boxtimes 
\sO_{\sZ^{\leq n}}).
\]

We also have a canonical adjunction morphism
$\gamma_*^{\IndCoh}(\omega_{\Gr_{\bG_m}^{pos,\leq n}}) \to 
\omega_{\Gr_{\bG_m}}$ for $\gamma:\Gr_{\bG_m}^{pos,\leq n} \to 
\Gr_{\bG_m}^{\leq n}$ the (ind-closed) embedding.

Pushing forward along $i_n:\sZ^{\leq n} \to \sY_{\log}^{\leq n}$
and composing, we obtain a canonical morphism:
%
\[
\Gamma^{\IndCoh}(\omega_{\Gr_{\bG_m}^{pos,\leq n}}) \otimes  
i_{n,*}^{\IndCoh}(\sO_{\sZ^{\leq n}}) \to 
\act_*^{\IndCoh}(\omega_{\Gr_{\bG_m}}^{\leq n} \boxtimes 
i_{n,*}^{\IndCoh}(\sO_{\sZ^{\leq n}}) = 
\Oblv \Av_!^{\Gr_{\bG_m}^{\leq n},w} 
i_{n,*}(\sO_{\sZ^{\leq n}}).
\]

% check notation on Av_!. 
% And should we emphasize averaging for \Gr_{\bG_m}^{\leq n}?

By Proposition \ref{p:cc-log}, we have a canonical isomorphism:
%
\[
\Gamma^{\IndCoh}(\omega_{\Gr_{\bG_m}^{pos,\leq n}}) \simeq
\Gamma(\fL_n^+ \bA^1,\sO_{\fL_n^+ \bA^1}).
\]

\noindent Therefore, by adjunction, the above gives a morphism:
%
\[
\begin{gathered}
\Gamma(\fL_n^+ \bA^1,\sO_{\fL_n^+ \bA^1}) \to 
\ul{\Hom}_{\IndCoh^*(\sY_{\log}^{\leq n})}
\big(i_{n,*}(\sO_{\sZ^{\leq n}}),
\Oblv\Av_!^{\Gr_{\bG_m}^{\leq n},w}  i_{n,*}(\sO_{\sZ^{\leq n}})\big) = \\
\ul{\End}_{\IndCoh^*(\sY^{\leq n})}
\big(\Av_!^{\Gr_{\bG_m}^{\leq n},w}  i_{n,*}(\sO_{\sZ^{\leq n}})\big).
\end{gathered}
\]

\noindent By \emph{loc. cit}., this is a morphism
of algebras. 

For $\omega \in t^{-n}k[[t]]dt/k[[t]]dt = (k[[t]]/t^n)^{\vee}$,
there is a corresponding linear function on $\fL_n^+ \bA^1$;
its image under the above map is by definition $\vph_{\omega}$ in 
\eqref{eq:fns}. As the above map extended to the symmetric
algebra on this vector space, we have verified that
the operators $\vph_{\omega}$ commute.

\begin{rem}
\source{tates-thesis-de-rham-setting}
\notready

More evocatively, but a little less rigorously, we have algebra maps:
%
\[
\Gr_{\bG_m}^{pos,\leq n} \to \Maps_{/\sY^{\leq n}}(\sZ^{\leq n},\sZ^{\leq n}) \to 
\ul{\End}_{\IndCoh^*(\sY^{\leq n})}(\sF_n)
\]

\noindent where the first map is given by the action and the second 
map sends a map $f:\sZ^{\leq n} \to \sZ^{\leq n}$ over $\sY^{\leq n}$ 
to the map: 
%
\[
\sF_n = \pi_{n,*}^{\IndCoh}(\sO_{\sZ^{\leq n}}) 
\to \pi_{n,*}^{\IndCoh}f_*^{\IndCoh}(\sO_{\sZ^{\leq n}}) = 
\pi_{n,*}^{\IndCoh}(\sO_{\sZ^{\leq n}}) = \sF_n
\]

\noindent where $\pi_n$ is the ind-proper morphism $\sZ^{\leq n} \to \sY^{\leq n}$ 
and the first map comes by functoriality from the canonical\footnote{This
map is tautologically equivalent via $\Psi$ to the adjunction morphism
$\sO_{\sZ} \to f_*^{\IndCoh}(\sO_{\sZ^{\leq n}}) \in \QCoh(\sZ^{\leq n})$.}
map $\sO_{\sZ^{\leq n}} \to f_*^{\IndCoh}(\sO_{\sZ^{\leq n}})$.

We then obtain an algebra map:
%
\[
\Gamma(\fL_n^+ \bA^1,\sO_{\fL_n^+ \bA^1}) =
\Gamma^{\IndCoh}(\omega_{\Gr_{\bG_m}^{pos,\leq n}}) \to 
\ul{\End}_{\IndCoh^*(\sY^{\leq n})}(\sF_n)
\]

\noindent from \eqref{eq:integration}. The detailed construction articulates
this idea explicitly.

\end{rem}

\subsection{Action of vector fields}

We now define the map $f \mapsto \xi_f$ from
\eqref{eq:vfs}.

\subsubsection{}

We begin with an elementary observation
involving $C = \Spec(k[x,y]/xy)$.
Let $i:\bA_y^1 \into C$ denote the 
embedding of the $y$-axis.

The map:
%
\[
x\cdot-:\sO_C \to \sO_C
\]

\noindent obviously factors through a map: 
%
\[
i_*\sO_{\bA_y^1} \to \sO_C
\]

\noindent that we also denote $y \cdot -$.

\subsubsection{}\label{ss:c-ind}
\source{tates-thesis-de-rham-setting}

We now vary to above to incorporate $\bG_m$-actions.

Consider $\bG_m \times \bG_m$ acting on
$C$ with the first $\bG_m$ scaling the $x$-axis
and the second $\bG_m$ scaling the $y$-axis.

The map $y\cdot-$ evidently has bidegree
$(0,1)$.
Therefore, on the stack $C/\bG_m\times\bG_m$, we obtain
the following. 

By definition, there are tautological line bundles 
$\sL_1$ and $\sL_2$ on $C/\bG_m\times \bG_m$
equipped with sections $x \in \sL_1$ and 
$y \in \sL_2$ such that
$xy = 0$ as a section of $\sL_1 \otimes \sL_2$. 

From the above, we obtain a canonical map:
%
\begin{equation}\label{eq:sections-ann}
\source{tates-thesis-de-rham-setting}
y\cdot -: i_*\sO_{0/\bG_m \times \bA_y^1/\bG_m} \to \sL_2.
\end{equation}

\subsubsection{}

Before applying the above, we introduce some more
notation.

Let $\Gr_{\bG_m}^{neg} \subset \Gr_{\bG_m}$ denote
the image of $\Gr_{\bG_m}^{pos}$ under the inversion
map $\Gr_{\bG_m} \isom \Gr_{\bG_m}$. 
We use similar notation for $\Gr_{\bG_m}^{neg,\leq n}$.

Note that $\AJ^{-1}:\widehat{\sD} \to \Gr_{\bG_m}$
maps into $\Gr_{\bG_m}^{neg}$. Similarly,
$\AJ^{-1}$ maps $\sD^{\leq n}$ into 
$\Gr_{\bG_m}^{neg,\leq n}$.

We let $\act_{neg}:\Gr_{\bG_m}^{neg} \times \sY_{\log}
\to \sY_{\log}$ denote the action map.
We use evident variants of this as well; crucially,
we also let $\act_{neg}$ denote the composition:
%
\[
\sD^{\leq n} \times \sZ^{\leq n} 
\xar{\AJ^{-1} \times i_n} 
\Gr_{\bG_m}^{neg} \times \sY_{\log} \xar{\act_{neg}}
\sY_{\log}.
\] 

\subsubsection{}\label{ss:act-inv}
\source{tates-thesis-de-rham-setting}

Now recall from Proposition \ref{p:z-flat-axes-var}
that we have a canonical map:
%
\[
\sD^{\leq n} \times \sZ^{\leq n} \to 
C/\bG_m \times \bG_m.
\]

In the notation of \S \ref{ss:c-ind},
the line bundle $\sL_2$ pulls
back to 
$\omega_{\sD^{\leq n}} \boxtimes \sO_{\sZ^{\leq n}}$  
by construction. The pullback of its section $y$ 
corresponds to the canonical map:
%
\[
\sZ^{\leq n} \to \LocSys_{\bG_m,\log}^{\leq n}
\xar{\Pol}
\fL^{pol,\leq n} \bA^1 dt = 
\Gamma(\sD^{\leq n},\omega_{\sD^{\leq n}})  
\] 

\noindent (with the last term really meaning the
scheme attached to this vector space). We sometimes
denote this section as $\omega^{univ}$ in what
follows.

We let 
$\act_{neg}^{-1}(\sZ^{\leq n}) \subset 
\sD^{\leq n} \times \sZ^{\leq n}$ denote the fiber
product:
%
\[
(\sD^{\leq n} \times \sZ^{\leq n})
\underset{C/\bG_m\times\bG_m}{\times}
(0/\bG_m \times \bA_y^1/\bG_m).
\]

\noindent (See \S \ref{ss:act-neg-rems} for
an explanation of the notation.) 
By Proposition \ref{p:z-flat-axes-var}, the
derived fiber product is classical.\footnote{This
fact is psychologically convenient, but not literally
necessary in what follows. I.e., a straightforward
restructuring of the discussion that follows could avoid 
directly appealing to this fact.} We let
$\alpha$ denote the embedding of
$\act_{neg}^{-1}(\sZ^{\leq n})$ into 
$\sD^{\leq n} \times \sZ^{\leq n}$.

Pulling back \eqref{eq:sections-ann}, we obtain 
a canonical map:
%
\begin{equation}\label{eq:supp}
\source{tates-thesis-de-rham-setting}
\alpha_*^{\IndCoh}
(\sO_{\act_{neg}^{-1}(\sZ^{\leq n})}) \to 
\omega_{\sD^{\leq n}} \boxtimes \sO_{\sZ^{\leq n}}
\in \IndCoh^*(\sD^{\leq n} \times \sZ^{\leq n})^{\heart}
\simeq \QCoh(\sD^{\leq n} \times \sZ^{\leq n})^{\heart}.
\end{equation}

\noindent (As the notation indicates, the
superscript $\IndCoh$ on $\alpha_*$ is a matter
of perspective: it is convenient for us for later 
purposes to view this morphism occurring in $\IndCoh^*$
as opposed to $\QCoh$.)

\begin{rem}
\source{tates-thesis-de-rham-setting}
\notready

In other words, the section $\omega^{univ}$
is scheme-theoretically supported on 
$\act_{neg}^{-1}(\sZ^{\leq n})$.

\end{rem}

\subsubsection{}\label{ss:act-neg-rems}
\source{tates-thesis-de-rham-setting}

We now claim that $\act_{neg}^{-1}(\sZ^{\leq n})$
is the \emph{classical} (not derived!) fiber
product:
%
\begin{equation}\label{eq:act-neg-rems}
\source{tates-thesis-de-rham-setting}
\act_{neg}^{-1}(\sZ^{\leq n}) = 
\Big(
(\sD^{\leq n} \times \sZ^{\leq n})
\underset{\sY_{\log}^{\leq n}}{\times}
\sZ^{\leq n} 
\Big)^{cl} \subset \sD^{\leq n} \times \sZ^{\leq n}.
\end{equation}

\noindent Here the morphism 
$\sD^{\leq n} \times \sZ^{\leq n} \to \sY_{\log}^{\leq n}$
is $\act_{neg}$.

Indeed, we have seen that 
$\act_{neg}^{-1}(\sZ^{\leq n})$
is a classical stack, so it suffices to check
the above identity on $A$-points for classical
commutative rings $A$, i.e., we are free to 
manipulate our $A$-points in the most naive way.

To calculate $\act_{neg}^{-1}(\sZ^{\leq n})$,
we need to recall from 
Proposition \ref{p:z-flat-axes-var} that
in the notation of \S \ref{ss:c-ind}, the
pullback of the line bundle $\sL$ comes from
the evident universal line bundle
on $\widehat{\sD} \times \sZ$; i.e.,
its fiber at a point 
$(\tau,(\sL,\nabla,s)) \in \widehat{\sD} \times \sZ$ 
is $\sL|_{\tau}$. Its canonical section,
denoted $x$ in \S \ref{ss:c-ind}, corresponds
to $s|_{\tau} \in \sL|_{\tau}$. 

Therefore, as a classical prestack,
$\act_{neg}^{-1}(\sZ^{\leq n}) \subset 
\sD^{\leq n} \times \sZ^{\leq n}$
corresponds to those data
$(\tau,(\sL,\nabla,s))$ with $s \in \sL(-\tau)$.

Clearly this description exactly matches
\eqref{eq:act-neg-rems}.

\begin{rem}
\source{tates-thesis-de-rham-setting}
\notready

The derived version of the
fiber product \eqref{eq:act-neg-rems} is easily
seen to differ from 
$\act_{neg}^{-1}(\sZ^{\leq n})$, so our notation
is a bit abusive from our perspective, which
emphasizes derived geometry. (However, we do
not feel so bad about this as we have provided
another derivedly good definition of the space.)

\end{rem}

\subsubsection{}\label{ss:can-constr}
\source{tates-thesis-de-rham-setting}

We now construct a canonical map:
%
\[
\on{can}:
i_{n,*}^{\IndCoh}\sO_{\sZ^{\leq n}} \to 
\act_{neg,*}^{\IndCoh}
(\omega_{\sD^{\leq n}} \boxtimes 
i_{n,*}\sO_{\sZ^{\leq n}}) \in 
\IndCoh^*(\sY_{\log}^{\leq n})
\]

\noindent as follows. 

By the above, we have a commutative diagram:
%
\[
\xymatrix{
\act_{neg}^{-1}(\sZ^{\leq n}) \ar[d]^{\act_{neg}^{\prime}} 
\ar[rr]^{\alpha} && 
\sD^{\leq n} \times \sZ^{\leq n} \ar[d]^{\act_{neg}} \\ 
\sZ^{\leq n} \ar[rr]^{i_n} && \sY_{\log}^{\leq n}.
}
\]

Now applying $i_{n,*}^{\IndCoh}$
to the evident counit map, we obtain:
%
\[
\begin{gathered}
i_{n,*}^{\IndCoh}\sO_{\sZ^{\leq n}} \to 
i_{n,*}^{\IndCoh} \act_{neg,*}^{\prime,\IndCoh}
\act_{neg}^{\prime,*,\IndCoh}(\sO_{\sZ^{\leq n}}) = 
i_{n,*}^{\IndCoh} \act_{neg,*}^{\prime,\IndCoh}
\sO_{\act_{neg}^{-1}(\sZ^{\leq n})} = \\
\act_{neg,*}^{\IndCoh} \alpha_*^{\IndCoh}
\sO_{\act_{neg}^{-1}(\sZ^{\leq n})}.
\end{gathered}
\]

\noindent We then applying $\act_{neg,*}^{\IndCoh}$
to \eqref{eq:supp} to obtain a map:
%
\[
\act_{neg,*}^{\IndCoh}\alpha_*^{\IndCoh}
(\sO_{\act_{neg}^{-1}(\sZ^{\leq n})}) \to 
\act_{neg,*}^{\IndCoh}
(\omega_{\sD^{\leq n}} \boxtimes \sO_{\sZ^{\leq n}}).
\]

\noindent Composing these two maps, we obtain
the desired map $\on{can}$.

\subsubsection{}

More generally, suppose $f \in k[[t]]/t^n$.

There is an associated action map 
$f\cdot -:\omega_{\sD^{\leq n}} \to 
\omega_{\sD^{\leq n}} \in \QCoh(\sD^{\leq n})$.
By functoriality, this gives rise to a morphism:
%
\[
f\cdot-:\act_{neg,*}^{\IndCoh}
(\omega_{\sD^{\leq n}} \boxtimes \sO_{\sZ^{\leq n}})
\to \act_{neg,*}^{\IndCoh}
(\omega_{\sD^{\leq n}} \boxtimes \sO_{\sZ^{\leq n}}).
\]

We then define: 
%
\[
\on{can}_f:
i_{n,*}^{\IndCoh}\sO_{\sZ^{\leq n}} \to 
\act_{neg,*}^{\IndCoh}
(\omega_{\sD^{\leq n}} \boxtimes 
i_{n,*}\sO_{\sZ^{\leq n}})
\]

\noindent as the composition of $\on{can} (= \on{can}_1)$
with $f\cdot-$.

\subsubsection{}

We now conclude the construction. 
Fix $f$ as above.

By adjunction,
there is a canonical morphism
$\AJ_*^{-1,\IndCoh}(\omega_{\sD^{\leq n}}) \to 
\omega_{\Gr_{\bG_m}^{\leq n}}$. 

Therefore, $\on{can}_f$ gives rise to a morphism:
%
\[
i_{n,*}^{\IndCoh}\sO_{\sZ^{\leq n}} \to 
\act_{neg,*}^{\IndCoh}
(\omega_{\Gr_{\bG_m}^{\leq n}} \boxtimes 
i_{n,*}\sO_{\sZ^{\leq n}}) = 
\Oblv\Av_!^{\Gr_{\bG_m}^{\leq n},w} 
(i_{n,*}\sO_{\sZ^{\leq n}}).
\]

\noindent By adjunction, we obtain the desired
morphism:
%
\[
\xi_f:\sF_n \to \sF_n.
\]

\subsubsection{}

Next, we show that the endomorphisms
$\xi_f$ of $\sF_n$ mutually commute.

By construction,\footnote{See \S \ref{ss:sigma-comp} for a much more general
setup.}
the assertion immediately reduces to the following one.

\begin{lem}
\source{tates-thesis-de-rham-setting}
\notready

The section:
%
\[
p_{13}^*(\omega^{univ}) \otimes p_{23}^*(\omega^{univ}) \in 
p_1^*(\omega_{\sD^{\leq n}}) \otimes p_2^*(\omega_{\sD^{\leq n}}) =
\omega_{\sD^{\leq n}} \boxtimes \omega_{\sD^{\leq n}} 
\boxtimes \sO_{\sZ^{\leq n}}
\]

\noindent is $\bZ/2$-equivariant (for the natural action
permuting the first two factors).

\end{lem}

\begin{proof}

There is a canonical morphism:
%
\[
\sD^{\leq n} \times \sD^{\leq n} \times \fL_n^+ \bA^1/\bG_m
\]

\noindent given as follows. To a point of the left hand side 
defined by
$\tau_1 \in \sD^{\leq n}$, $\tau_2 \in \sD^{\leq n}$,
a line bundle $\sL$ on $\sD^{\leq n}$ and a section 
$s \in \sL$, our map assigns the line
$\sL|_{\tau_1} \otimes \sL|_{\tau_2}$ with
its section $s|_{\tau_1}\otimes s|_{\tau_2}$.

In fact, this map factors through a map:
%
\[
\Sym^2(\sD^{\leq n}) \times \fL_n^+ \bA^1/\bG_m.
\]

\noindent Indeed, for a point $(D,(\sL,s))$ of the left hand
side, we may regard $D$ as a finite subscheme of $\sL$.
The map then sends the datum to 
$(\Lambda^2(s) \in \Lambda^2(\sL|_D)) \in \bA^1/\bG_m$.
I.e., this is evident by the usual norm construction.

Unwinding the constructions, this implies the claim.

\end{proof}

\begin{rem}
\source{tates-thesis-de-rham-setting}
\notready

We could have organized the discussion differently. 
Instead, we could have generalized our
construction of $\xi_f$ with $\Gr_{\bG_m}^{neg,\leq n}$ replacing
$\sD^{\leq n}$ via a suitable use of norms (as in the lemma
above). Then we would immediately deduce commutativity of the
operators $\xi_f$ by the same argument and for the operators
$\vph_{\omega}$.

\end{rem}

\subsection{Uncertainty}

We now check the remaining Weyl algebra relation
\eqref{eq:uncertainty}.

\subsubsection{}\label{ss:gpd}
\source{tates-thesis-de-rham-setting}

First, let us put the above constructions a common framework to allow us to 
compute the relevant compositions.

Let $\sH^{\leq n} \coloneqq \sZ^{\leq n} \times_{\sY^{\leq n}} \sZ^{\leq n}$.
This stack is a groupoid over $\sZ^{\leq n}$ with its projections
$p_1,p_2:\sH^{\leq n} \rightrightarrows \sZ^{\leq n}$ ind-proper.

By adjunction and base-change, morphisms $\sF_n \to \sF_n$ are equivalent to
sections of $p_1^!(\sO_{\sZ^{\leq n}})$. For\footnote{For $V \in \Vect$,
$v\in V$ means $v \in \Omega^{\infty} V$, i.e., we have a point of the
underlying $\infty$-groupoid. If we concretely model $V$ as a cochain
complex, we obtain such data from cycles in degree $0$.}
$\sigma \in \Gamma^{\IndCoh}(\sH^{\leq n},p_1^!(\sO_{\sZ^{\leq n}}))$, 
we let $\psi_{\sigma}:\sF_n \to \sF_n$ denote the corresponding morphism.

In particular, suppose we are given a prestack $S$ with an ind-proper morphism
$\eta:S \to \sH^{\leq n}$. Given a section of 
$\eta^!p_1^!(\sO_{\sZ^{\leq n}})$, we obtain a section of 
$p_1^!(\sO_{\sZ^{\leq n}})$ by functoriality, hence a morphism
as above.

\begin{example}
\source{tates-thesis-de-rham-setting}
\notready\label{e:fns}
\source{tates-thesis-de-rham-setting}

Let $S = \sD^{\leq n} \times \sZ^{\leq n}$ and let
$\eta:S \to \sH^{\leq n}$
be the morphism: 
%
\[
(p_2,\act\circ (\AJ \times \id)):
(\tau,(\sL,\nabla,s) \mapsto ((\sL,\nabla,s),(\sL(\tau),\nabla,s)) \in 
\sZ^{\leq n} \underset{\sY^{\leq n}}{\times} \sZ^{\leq n} = 
\sH^{\leq n}.
\]

For $\omega \in (k[[t]]/t^n)^{\vee} \overset{Prop. \ref{p:cc-log}}{=} 
\Gamma^{\IndCoh}(\sD^{\leq n},\omega_{\sD^{\leq n}})$, we
obtain the evident section $p_1^*(\omega)$ of:
%
\[
\omega_{\sD^{\leq n}} \boxtimes \sO_{\sZ^{\leq n}} = 
p_2^! \sO_{\sZ^{\leq n}}  = 
\eta^!p_1^!(\sO_{\sZ^{\leq n}}).
\]

\noindent The corresponding map $\sF_n \to \sF_n$ is
\eqref{eq:fns}.

\end{example}

\begin{example}
\source{tates-thesis-de-rham-setting}
\notready\label{e:vfs}
\source{tates-thesis-de-rham-setting}

In the notation of \S \ref{ss:act-inv}, let 
$S = \act_{neg}^{-1}(\sZ^{\leq n}) \subset \sD^{\leq n} \times \sZ^{\leq n}$.
Let $\eta$ be the evident morphism defined using \eqref{eq:act-neg-rems}.\footnote{In the notation of \S \ref{ss:can-constr}, this 
morphism is given by $(p_1\alpha,\act_{neg}^{\prime})$.}
The construction of $\on{can}$ (resp. $\on{can}_f$) amounted to showing
that the section $\omega^{univ}$ (hence $p_1^*(f)\cdot \omega^{univ}$) 
comes from a canonical section of
$\eta^!p_1^!(\sO_{\sZ^{\leq n}})$ (i.e., these sections are scheme-theoretically
supported on $\act_{neg}^{-1}(\sZ^{\leq n}) $). By construction,
the corresponding map $\sF_n \to \sF_n$ is \eqref{eq:vfs}.

\end{example}

\begin{rem}
\source{tates-thesis-de-rham-setting}
\notready\label{r:h-cl}
\source{tates-thesis-de-rham-setting}

Although 
$\sH^{\leq n}$ is not classical, our constructions
here are not sensitive to this. The reason is that
we have a closed embedding $\iota:\sH^{\leq n} \into 
\Gr_{\bG_m}^{\leq n} \times \sZ^{\leq n}$ 
(cf. \S \ref{ss:h-embed} below). As 
$\Gr_{\bG_m}^{\leq n}$ is ind-finite and 
$\sZ^{\leq n}$ is classical, 
$\Gamma(p_1^!(\sO_{\sZ^{\leq n}})) = \Gamma(\iota^!(\omega_{\Gr_{\bG_m}^{\leq n}}
\boxtimes \sO_{\sZ^{\leq n}}))$ is a coconnective complex,
and its $H^0$ is the same for $\sH^{\leq n}$ and for its
underlying classical stack (both identify with 
sections of $\omega_{\Gr_{\bG_m}^{\leq n}}
\boxtimes \sO_{\sZ^{\leq n}}$ scheme-theoretically supported
on $\sH^{\leq n,cl}$).

\end{rem}

\subsubsection{}\label{ss:sigma-comp}
\source{tates-thesis-de-rham-setting}

We now spell out how to concretely compose morphisms of the above type.

We have three morphisms:
%
\[
p_{12},p_{23},p_{13}:
\sH^{\leq n} \underset{\sZ^{\leq n}}{\times} \sH^{\leq n} = 
\sZ^{\leq n} \underset{\sY^{\leq n}}{\times}
\sZ^{\leq n} \underset{\sY^{\leq n}}{\times} \sZ^{\leq n}
\to 
\sZ^{\leq n} \underset{\sY^{\leq n}}{\times} \sZ^{\leq n} =
\sH^{\leq n}.
\]

\noindent We remark that $p_{13}$ corresponds to the multiplication
on the groupoid.

%Let $m:\sH^{\leq n} \times_{\sZ^{\leq n}} \sH^{\leq n} \to \sH^{\leq n}$ be the 
%multiplication for the groupoid (so the maps 
%$\sH^{\leq n} \to \sZ^{\leq n}$ are $p_2$ and $p_1$ respectively). 
We begin by claiming that there is a canonical 
isomorphism: % this is an iso where???

\begin{equation}\label{eq:p13-!}
\source{tates-thesis-de-rham-setting}
p_{13}^!p_1^!(\sO_{\sZ^{\leq n}}) \simeq 
p_{12}^*p_1^!(\sO_{\sZ^{\leq n}}) \otimes 
p_{23}^*p_1^!(\sO_{\sZ^{\leq n}}). 
\end{equation}

\noindent To see this, first observe that
$p_1p_{13} = p_1p_{12}$, so:
%
\[
p_{13}^!p_1^!(\sO_{\sZ^{\leq n}}) = 
p_{12}^!p_1^!(\sO_{\sZ^{\leq n}}).
\]

Next, we have evident isomorphisms:
%
\[
\begin{gathered}
p_{12}^!p_1^! \sO_{\sZ^{\leq n}} = 
p_{12}^!(p_1^! \sO_{\sZ^{\leq n}}
\underset{\sO_{\sH^{\leq n}}}{\otimes} 
\sO_{\sH^{\leq n}}) = 
p_{12}^*p_1^! \sO_{\sZ^{\leq n}}
\underset{\sH^{\leq n} \underset{\sZ^{\leq n}} \sH^{\leq n} }
{\otimes} 
p_{12}^!\sO_{\sH^{\leq n}}.
\end{gathered}
\]

\noindent Then the Cartesian diagram:
%
\[
\xymatrix{
\sH^{\leq n} \underset{\sZ^{\leq n}}{\times} \sH^{\leq n} \ar[rr]^{p_{12}} 
\ar[d]_{p_{23}} && \sH^{\leq n} \ar[rr]^{p_1} \ar[d]^{p_2} && 
\sZ^{\leq n} \\
\sH^{\leq n} \ar[rr]^{p_1} && \sZ^{\leq n}
}
\]

\noindent yields an isomorphism: % reference?
%
\[
p_{12}^!\sO_{\sH^{\leq n}} = 
p_{12}^!p_2^*\sO_{\sZ^{\leq n}} = 
p_{23}^*p_1^!\sO_{\sZ^{\leq n}}.
\]

\noindent Combining the above isomorphims yields \eqref{eq:p13-!}.

Therefore, given $\sigma_1, \sigma_2$ sections of
$p_1^!(\sO_{\sZ^{\leq n}})$, we obtain a section:
%
\[
p_{12}^*(\sigma_1) \otimes p_{23}^*(\sigma_2)
\]

\noindent of $p_{13}^!p_1^!(\sO_{\sZ^{\leq n}})$
(using \eqref{eq:p13-!}). As $p_{13}$ is ind-proper,
we obtain an induced section $\sigma_1 \star \sigma_2$ 
of $p_1^!(\sO_{\sZ^{\leq n}})$. 

Then unwinding constructions, we find:
%
\[
\psi_{\sigma_2} \psi_{\sigma_1} = \psi_{\sigma_1 \star \sigma_2}.
\]

% ordering?

\subsubsection{}\label{ss:h-embed}
\source{tates-thesis-de-rham-setting}

To analyze the composition above, 
it is convenient to embed our spaces as follows.

First, define a morphism:
%
\[
\iota:\sH^{\leq n} \into \Gr_{\bG_m}^{\leq n} \times \sZ^{\leq n}
\]

\noindent as:
%
\[
\sH^{\leq n} = 
\sZ^{\leq n} \underset{\sY^{\leq n}}{\times} \sZ^{\leq n} 
\into 
\sZ^{\leq n} \underset{\sY^{\leq n}}{\times} \sY_{\log}^{\leq n}
\overset{\simeq}{\leftarrow} \Gr_{\bG_m}^{\leq n} \times \sZ^{\leq n} 
\]

\noindent where the last arrow is $(p_2,\act)$. Note that
the composition of $\iota$ with the projection to $\sZ^{\leq n}$
is the projection $p_1:\sH^{\leq n} \to \sZ^{\leq n}$.

Next, we define a similar morphism:
%
\[
\widetilde{\iota}:
\sH^{\leq n} \underset{\sZ^{\leq n}}{\times}
\sH^{\leq n} \into 
\Gr_{\bG_m}^{\leq n} \times \Gr_{\bG_m}^{\leq n} \times \sZ^{\leq n}
\]

\noindent as:
%
\[
\sH^{\leq n} \underset{\sZ^{\leq n}}{\times}
\sH^{\leq n} = 
\sZ^{\leq n} \underset{\sY^{\leq n}}{\times} \sZ^{\leq n} 
\underset{\sY^{\leq n}}{\times} \sZ^{\leq n} 
\into 
\sZ^{\leq n} \underset{\sY^{\leq n}}{\times} \sY_{\log}^{\leq n}
\underset{\sY^{\leq n}}{\times} \sY_{\log}^{\leq n}
\overset{\simeq}{\leftarrow} 
\Gr_{\bG_m}^{\leq n} \times \Gr_{\bG_m}^{\leq n} \times \sZ^{\leq n} 
\]

\noindent where this time the last arrow is 
$(p_3,\act \circ p_{23}, \act \circ (\id \times \act))$. 

As in Remark \ref{r:h-cl}, a section of 
$\Gamma^{\IndCoh}(\sH^{\leq n},p_1^!(\sO_{\sZ^{\leq n}}))$
is the same as a section of 
$\omega_{\Gr_{\bG_m}^{\leq n}} \boxtimes \sO_{\sZ^{\leq n}}$
scheme-theoretically supported on $\sH^{\leq n}$ 
(equivalently, on the underlying classical stack).
The same applies for 
$\widetilde{\iota}$: a section
of $p_{13}^!p_1^!(\sO_{\sZ^{\leq n}})$ is the same
as a section of:
%
\[
\omega_{\Gr_{\bG_m}^{\leq n}} \boxtimes \omega_{\Gr_{\bG_m}^{\leq n}} \boxtimes \sO_{\sZ^{\leq n}}
\]

\noindent scheme-theoretically supported on 
$\sH^{\leq n} \times_{\sZ^{\leq n}} \sH^{\leq n}$ (which is
equivalent to saying at the 
classical level).

\subsubsection{}\label{sss:recipe}
\source{tates-thesis-de-rham-setting}

Now fix $\omega \in (k[[t]]/t^n)^{\vee}$ and $f \in k[[t]]/t^n$.
In effect, \S \ref{ss:sigma-comp} constructed two sections of
$p_{13}^!p_1^!(\sO_{\sZ^{\leq n}})$, which
induce the compositions 
$\xi_f\vph_{\omega}$ and $\vph_{\omega}\xi_f$ (after 
pushforward along $p_{13}$).

Let us explicitly describe the resulting sections of 
$\omega_{\Gr_{\bG_m}^{\leq n}} \boxtimes \omega_{\Gr_{\bG_m}^{\leq n}} 
\boxtimes \sO_{\sZ^{\leq n}}$ under the
above dictionary.

For $\vph_{\omega}\xi_f$, the corresponding section
is supported on:
%
\[
\sD^{\leq n} \times \sD^{\leq n} \times \sZ^{\leq n}
\overset{\AJ \times \AJ^{-1} \times \id}{\into}
\Gr_{\bG_m}^{\leq n} \times \Gr_{\bG_m}^{\leq n} 
\times \sZ^{\leq n}.
\]

\noindent The corresponding section of:
%
\[
\omega_{\sD^{\leq n}} \boxtimes 
\omega_{\sD^{\leq n}} \boxtimes \sO_{\sZ^{\leq n}}
\]

\noindent is $\omega \boxtimes p_2^*(f)\cdot\omega^{univ}$,
with notation as in Examples \ref{e:fns} and \ref{e:vfs}.\footnote{As
we have shown, this section is scheme-theoretically supported
on $\sD^{\leq n} \boxtimes \act_{neg}^{-1}(\sZ^{\leq n})$,
hence on $(\sH^{\leq n} \times_{\sZ^{\leq n}}\sH^{\leq n})^{cl}$.} It is convenient to 
rewrite this section as:
%
\[
p_2^*(f) \cdot (p_1^*(\omega) \otimes p_{23}^*(\omega^{univ})) \in 
p_1^*(\omega_{\sD^{\leq n}}) \otimes 
p_{23}^*(\omega_{\sD^{\leq n}} \boxtimes \sO_{\sZ^{\leq n}}).
\]

For $\xi_f\vph_{\omega}$, the corresponding section
is supported on:
%
\[
\sD^{\leq n} \times \sD^{\leq n} \times \sZ^{\leq n}
\xar{\AJ^{-1} \times \AJ \times \id}
\Gr_{\bG_m}^{\leq n} \times \Gr_{\bG_m}^{\leq n} 
\times \sZ^{\leq n}.
\]

\noindent The corresponding section of:
%
\[
\omega_{\sD^{\leq n}} \boxtimes 
\omega_{\sD^{\leq n}} \boxtimes \sO_{\sZ^{\leq n}}
\]

\noindent is:\footnote{For clarity, we again
highlight that we have shown
this section is in fact supported on 
$(\sH^{\leq n} \times_{\sZ^{\leq n}}\sH^{\leq n})^{cl}$, as
it should be. Indeed, at the classical level, our previous
analysis shows that it is
supported on 
$(\act_{neg}\circ (\id\times\act))^{-1} (\sZ^{\leq n}) 
\subset  (\sH^{\leq n} \times_{\sZ^{\leq n}}\sH^{\leq n})^{cl}$.}
%
\[
\begin{gathered}
p_1^*(f) \cdot 
(p_2^*(\omega) \otimes (p_1\times \act)^*(\omega^{univ})) \in 
p_2^*(\omega_{\sD^{\leq n}}) \otimes 
(p_1\times \act)^*(\omega_{\sD^{\leq n}} \boxtimes \sO_{\sZ^{\leq n}}) = \\
\omega_{\sD^{\leq n}} \boxtimes \omega_{\sD^{\leq n}} 
\boxtimes \sO_{\sZ^{\leq n}}.
\end{gathered}
\]

Now note that we have a commutative diagram:
%
\[
\xymatrix{
\sH^{\leq n} \underset{\sZ^{\leq n}}{\times} \sH^{\leq n} 
\ar[rr]^{\widetilde{\iota}} \ar[d]_{p_{13}} &&
\Gr_{\bG_m}^{\leq n} \times \Gr_{\bG_m}^{\leq n} \times \sZ^{\leq n} 
\ar[d]^{m\times \id} \\
\sH^{\leq n} \ar[rr]^{\iota} && \Gr_{\bG_m}^{\leq n} \times \sZ^{\leq n}
}
\]

\noindent for $m$ the multiplication on $\Gr_{\bG_m}^{\leq n}$.
As this multiplication is commutative, we see that
$\xi_f\vph_{\omega}$ could as well have been defined
by the section as above obtained by transposing
the first two coordinates, i.e.:
%
\[
\begin{gathered}
p_2^*(f) \cdot 
(p_1^*(\omega) \otimes (p_2,\act \circ p_{13})^*(\omega^{univ})).
\end{gathered}
\]

\noindent This section has the advantage of
sharing its support with our previous one for 
$\vph_{\omega}\xi_f$.

\subsubsection{}

We can now conclude the argument. We use the next elementary lemmas.

\begin{lem}
\source{tates-thesis-de-rham-setting}
\notready\label{l:cotr-1}
\source{tates-thesis-de-rham-setting}

Let: 
%
\[
\cotr \in 
\Gamma(\sD^{\leq n},\sO_{\sD^{\leq n}}) \otimes
\Gamma(\sD^{\leq n},\sO_{\sD^{\leq n}})^{\vee} = 
\Gamma^{\IndCoh}(\sD^{\leq n} \times \sD^{\leq n},
\sO_{\sD^{\leq n}} \boxtimes \omega_{\sD^{\leq n}})
\]

\noindent denote the canonical vector. 

We then have an equality:
%
\[
p_{23}^*(\omega^{univ}) = 
(p_2,\act \circ p_{13})^*(\omega^{univ}))
+p_{12}^*(\cotr)
\in 
\sO_{\sD^{\leq n}} \boxtimes 
\omega_{\sD^{\leq n}} \boxtimes \sO_{\sZ^{\leq n}}.
\]

\end{lem}

\begin{proof}

For convenience, we give the first (resp. second) factor
of our triple product as 
$\sD_{\tau_1}^{\leq n}$ (resp. $\sD_{\tau_2}^{\leq n}$), 
where $\tau_1,\tau_2$ denote the respective coordinates. 

For a prestack $S$, note that a section of 
$\omega_{\sD_{\tau_2}^{\leq n}} \boxtimes \sO_S$ is equivalent to
a map $S \to \fL^{pol,\leq n} \bA^1 dt$.

For $S = \sD_{\tau_1}^{\leq n} \times \sZ^{\leq n}$,
our sections above correspond to maps:
%
\[
\sD_{\tau_1}^{\leq n} \times \sZ^{\leq n} 
\xar{(\tau_1,(\sL,\nabla,s)) \mapsto ???}
\fL^{pol,\leq n} \bA^1 dt
\]

\noindent as follows:

\begin{align*}
p_{23}^*(\omega^{univ}) & \leftrightsquigarrow 
%\sD_{\tau_1}^{\leq n} \times \sZ^{\leq n} \xar{p_2} 
%\sZ^{\leq n} \xar{\Pol} \fL^{pol,\leq n} \bA^1 dt \\
(\tau_1,(\sL,\nabla,s)) \mapsto 
\Pol(\nabla) \\
(p_2,\act \circ p_{13})^*(\omega^{univ})) 
& \leftrightsquigarrow 
%\sD_{\tau_1}^{\leq n} \times \sZ^{\leq n} \xar{\act} 
%\sZ^{\leq n} \xar{\Pol} \fL^{pol,\leq n} \bA^1 dt \\
(\tau_1,(\sL,\nabla,s)) \mapsto 
\Pol((\sL(\tau_1),\nabla,s)) = \\ 
& \hspace{2cm} \Pol(\nabla)+d\log\AJ(\tau_1) \\
p_{12}^*(\cotr) & \leftrightsquigarrow 
%\sD_{\tau_1}^{\leq n} \times \sZ^{\leq n} \xar{p_1}
%\sD_{\tau_1}^{\leq n} \xar{d\log \AJ} \fL^{pol,\leq n} \bA^1 dt
(\tau_1,(\sL,\nabla,s)) \mapsto 
-d\log\AJ(\tau_1)
\end{align*}

\noindent where in the last two lines the map
$d\log\AJ$ was considered in
\S \ref{ss:dlog-aj}-\ref{ss:aj-trun}; the last line
follows from Lemma \ref{l:dlog-aj}.

We clearly obtain the assertion.

\end{proof}

We also use the following observations.

\begin{lem}
\source{tates-thesis-de-rham-setting}
\notready\label{l:cotr-2}
\source{tates-thesis-de-rham-setting}

The above section $\cotr$ of 
$\sO_{\sD^{\leq n}} \boxtimes \omega_{\sD^{\leq n}}$
is scheme-theoretically supported on the diagonally embedded
$\sD^{\leq n} \subset \sD^{\leq n} \times \sD^{\leq n}$.

Moreover, for $f \in k[[t]]/t^n$ and 
$\omega \in t^{-n}k[[t]]dt/k[[t]]dt$, the section:
%
\[
p_2^*(f) \cdot (p_1^*(\omega) \otimes \cotr) \in 
\omega_{\sD^{\leq n}} \boxtimes \omega_{\sD^{\leq n}} 
\]

\noindent maps to $\Res(f\omega) \in k$ under the
adjunction morphism:
%
\[
\Gamma^{\IndCoh}(\sD^{\leq n} \times \sD^{\leq n},
\omega_{\sD^{\leq n}} \boxtimes \omega_{\sD^{\leq n}} ) \to k.
\]

\end{lem}

\begin{proof}

Let $A$ be a commutative, classical, finite $k$-algebra.

As $\id_A:A \to A$ is a morphism of $A$-bimodules,
the corresponding element $\cotr \in A^{\vee} \otimes A \in 
A\otimes A\mod^{\heart}$ is scheme-theoretically 
supported on the diagonal 
$\Spec(A) \subset \Spec(A) \times \Spec(A)$.

Moreover, for $f \in A$ and $\omega \in A^{\vee}$, we
have an element $f \otimes \omega \in A \otimes A^{\vee}$.
Tensoring over $A\otimes A$ with $A^{\vee} \otimes A$,
we obtain an element:
%
\[
f\otimes \omega \cdot \cotr \in A^{\vee} \otimes A^{\vee}.
\]

\noindent It is easy to see that when we evaluate this
tensor on $1 \otimes 1 \in A \otimes A$, we obtain
$\omega(f) \in k$.

Taking $A = k[[t]]/t^n$, we obtain our claim.

\end{proof}

By our earlier discussion, the endomorphism 
$\xi_f \vph_{\omega}-\vph_{\omega}\xi_f$ of $\sF_n$
corresponds to the section:
%
\[
p_2^*(f) \cdot 
(p_1^*(\omega) \otimes (p_2,\act \circ p_{13})^*(\omega^{univ})) -
p_2^*(f) \cdot 
(p_1^*(\omega) \otimes p_{23}^*(\omega^{univ}))
\]

\noindent of $\omega_{\sD^{\leq n}} \boxtimes 
\omega_{\sD^{\leq n}} \boxtimes \sO_{\sZ^{\leq n}}$.\footnote{Of 
course we consider
$\sD^{\leq n} \times \sD^{\leq n} \times
\sZ^{\leq n}$ mapping to $\Gr_{\bG_m}^{\leq n} \times \sZ^{\leq n}$
via $(\tau_1,\tau_2,(\sL,\nabla,s)) \mapsto 
(\AJ(\tau_1)\cdot\AJ^{-1}(\tau_2),(\sL,\nabla,s))$.}

By Lemma \ref{l:cotr-1}, the above section coincides with:
%
\[
-p_2^*(f) \cdot p_{12}^*(p_1^*(\omega)\otimes\cotr).
\]

\noindent By Lemma \ref{l:cotr-2}, this section is supported
on the diagonally embedded:
%
\[
\sD^{\leq n} \times \sZ^{\leq n} 
\xar{\Delta \times \id} 
\sD^{\leq n} \times \sD^{\leq n} \times \sZ^{\leq n},
\]
 
\noindent which evidently
maps to $1 \times \sZ^{\leq n} \subset \Gr_{\bG_m}^{\leq n} \times \sZ^{\leq n}$
under $p_{13}$. Moreover, Lemma \ref{l:cotr-2} implies that the
pushforward of the above section to $\sZ^{\leq n} = 1 \times \sZ^{\leq n}$
is: 
%
\[
-\Res(f\omega) \in k \subset \Gamma(\sO_{\sZ^{\leq n}}) \subset 
\Gamma^{\IndCoh}(\Gr_{\bG_m}^{\leq n} \times \sZ^{\leq n}, 
\omega_{\Gr_{\bG_m}^{\leq n}} \boxtimes \sZ^{\leq n}).
\]

\noindent This amounts to \eqref{eq:uncertainty}, concluding
the argument.




%
%
%\section{Explicit analysis in the regular singular case}\label{s:rs}
%
%\subsection{}
